%% =========================================================================
\subsection{The Factor Tree and GenMixedMC}
\label{sec:factor-tree}
%% =========================================================================

The random variant's provability motivates a broader framework:
instead of fixing a single selection rule, consider the
\textbf{factor tree} $\mathcal{T}(m)$ rooted at starting point~$m$.
The root holds accumulator~$m$; each node with accumulator~$P$ has
one child $P \cdot p$ for every prime $p \mid P{+}1$.
The leftmost path (always choosing $\minFac$) recovers the standard EM
walk; every other path corresponds to a different selection.
All nodes of $\mathcal{T}(m)$ share the same structural properties:
each factor is prime, divides $P{+}1$, and is coprime to~$P$; the
accumulator is strictly increasing along every branch.

\begin{definition}[GenMixedMC]
\label{def:gen-mixed-mc}
$\mathrm{GenMixedMC}(m)$: for every prime~$q$, some node of
$\mathcal{T}(m)$ has $q$ dividing its Euclid number $P{+}1$.
Equivalently, every prime appears somewhere in the factor tree.
\end{definition}

\noindent
The standard Mullin Conjecture is the special case where the
leftmost path alone captures every prime; GenMixedMC relaxes this to
any path.  Formally, $\mathrm{MC} \Longrightarrow
\mathrm{GenMixedMC}(2)$, since the standard walk is one particular
path through $\mathcal{T}(2)$.

\begin{remark}[The Euclid--Mullin graph]
\label{rem:em-graph}
The factor tree $\mathcal{T}(m)$ is exactly the \emph{Euclid--Mullin
graph}~$G_m$ of Booker and Irvine~\cite{BoIr2016}.  In their
framework, vertices are accumulated products and edges are labelled by
primes dividing $P{+}1$; GenMixedMC$(m)$ asks whether every prime
labels some edge.  Computationally, every prime below~$10^9$ occurs as
an edge in~$G_1$ (hence in~$G_2$).

Booker~\cite{BoSa2012} proved that for any finite set~$P$ of primes,
there exists a path through $\mathcal{T}(\prod P)$ containing every
prime.  His proof uses three ingredients that connect directly to the
formalization's infrastructure:
\begin{enumerate}[nosep]
\item \emph{Squarefree smooth density}:
  the set~$S_q$ of residues mod~$q$ attained by squarefree
  $(q{-}1)$-smooth integers satisfies $|S_q| > (q{-}1)/2$
  (via Schnirelmann density $\geq 53/88$)---the same population
  equidistribution one would like to transfer to a single orbit
  (a transfer that no population statement supplies: Dead End~\#160).
\item \emph{Weil--Hasse bounds} for the curves $y^2 = x(x^2+a)$
  and $y^2 = x^6+a$---the same character-sum cancellation exploited
  by the function field analog (\S\ref{sec:function-field}).
\item \emph{Subgroup growth}: the subgroup of $(\ZZ/q\ZZ)^\times$
  generated by the reachable primes has index dividing~$6$---an
  analogue of SubgroupEscape (Theorem~\ref{thm:pre}), which guarantees
  full generation for the standard walk.
\end{enumerate}
Booker's proof works because choosing subsets at each step provides
enough freedom to steer residues mod~$q$.  The standard EM sequence
($\minFac$, always using the full product) has no such freedom: the
accumulator is deterministic, and the gap between population
equidistribution (Booker's Lemma~2) and orbit-specific
equidistribution (CME) is exactly the open hypothesis.
\end{remark}

\begin{remark}[Quantifier structure]
\label{rem:quantifier}
GenMixedMC$(m)$ asserts: for each prime~$q$, \emph{some} path
through $\mathcal{T}(m)$ reaches a node where $q$ divides the Euclid
number.  Different primes may require different paths.  The quantifier
order is $\forall q\;\exists\,\text{path}$, \textbf{not}
$\exists\,\text{path}\;\forall q$.  The stronger statement---a single
path capturing all primes simultaneously---remains open even for
almost all starting points.
\end{remark}

\paragraph{Lean encoding: MixedSelection.}
The formalization encodes paths through $\mathcal{T}(m)$ via a
\texttt{MixedSelection}
$\sigma : \NN \to \mathrm{Option}\;\NN$
(\lean{EM/Stochastic/MixedWalk.lean}{MixedSelection}):
$\mathrm{none}$ at step~$k$ means follow the leftmost child
($\minFac$); $\mathrm{some}\;p$ means branch to child~$p$.
A selection is \defname{valid}
(\lean{EM/Stochastic/MixedWalk.lean}{ValidMixedSelection})
if every $\mathrm{some}\;p$ step has $p$ prime and $p \mid P_\sigma(k){+}1$.
The accumulator along a path is
$P_\sigma(0) = m$,
$P_\sigma(n{+}1) = P_\sigma(n) \cdot f_\sigma(n)$
(\lean{EM/Stochastic/MixedWalk.lean}{mixedWalkProd}).

\paragraph{The reachable set: tree projection.}
Define $R_n(q, m)$ as the set of positions
$P_\sigma(n) \bmod q$ over all valid paths~$\sigma$
(\lean{EM/Stochastic/ReachableSets.lean}{reachableAt})---the
projection of depth-$n$ nodes of $\mathcal{T}(m)$ onto
$(\ZZ/q\ZZ)^\times$.  The cumulative set
$R_\infty = \bigcup_n R_n$ collects all residues attained by
any node.  Hitting is equivalent to reachability:
\[
  q \text{ appears in } \mathcal{T}(m)
  \;\Longleftrightarrow\;
  {-}1 \in R_\infty(q, m)
\]
(\lean{EM/Stochastic/ReachableSets.lean}{mixed\_hitting\_iff\_neg\_one\_reachable}
{\color{leanink}$\checkmark$}).
This reformulates GenMixedMC as a BFS reachability problem
in $(\ZZ/q\ZZ)^\times$.
The reachable set is monotone along walk prefixes:
any position reachable from a descendant $P_\sigma(n)$ is also
reachable from the original accumulator~$m$
(\lean{EM/Stochastic/TreeSieveDecay.lean}{reachableEver\_mono\_along\_walk}
{\color{leanink}$\checkmark$}),
via concatenation of walk segments.

\paragraph{Reachable set growth.}
The core growth mechanism: at any node with accumulator~$P$,
\emph{every} prime~$p$ dividing $P{+}1$ produces a child at position
$P \cdot p \bmod q$
(\lean{EM/Stochastic/ReachableSets.lean}{reachableAt\_from\_factor}
{\color{leanink}$\checkmark$}).
When $P{+}1$ is composite, it has $\geq 2$ prime divisors,
giving $\geq 2$ children
(\lean{EM/Stochastic/ReachableSets.lean}{reachable\_composite\_branch}
{\color{leanink}$\checkmark$}).
When these are distinct mod~$q$, $|R_{n+1}| > |R_n|$ along this branch.

\paragraph{The purification lemma.}
The probabilistic argument for the random variant requires
Diaconis--Shahshahani mixing theory and Borel--Cantelli, neither
available in Lean.  The formalization instead proves that the
nondeterministic framework captures everything:

\smallskip\noindent
\textbf{Purification Lemma}
(\lean{EM/Stochastic/RandomFactorMC.lean}{pure\_random\_mc\_iff\_mixed\_mc}
{\color{leanink}$\checkmark$}).
\emph{PureRandomMC$(q)$ $\Leftrightarrow$ MixedMC$(q)$.}

\smallskip\noindent
The proof constructs, from any path through $\mathcal{T}(m)$ (which
may follow the leftmost child at some steps), a
\emph{purified} path that always branches explicitly and follows the
exact same trajectory
(\lean{EM/Stochastic/RandomFactorMC.lean}{purify\_walk\_eq}
{\color{leanink}$\checkmark$}).
This works because $\minFac(P{+}1)$ is itself a prime dividing~$P{+}1$,
so replacing every ``follow leftmost'' with ``branch to
$\minFac(P{+}1)$'' produces a valid explicit path.
The distinction between ``$\minFac$-first'' and
``all-random'' is thus illusory at the nondeterministic level.
Standard~MC implies GenMixedMC$(2)$ (the leftmost path is one
particular branch,
\lean{EM/Stochastic/RandomFactorMC.lean}{standard\_mc\_implies\_pure\_random}
{\color{leanink}$\checkmark$}).

\paragraph{Variant hierarchy.}
The factor tree framework unifies the prior variants in a strict
hierarchy:
\[
  \text{deterministic } (\text{leftmost path})
  \;\subset\;
  \text{two-point } (\S\ref{sec:minfac-secondminfac})
  \;\subset\;
  \mathcal{T}(m)
  \;\subset\;
  \text{full random.}
\]
The two-point walk (\S\ref{sec:minfac-secondminfac}) restricts
branching to the two smallest factors; the full random walk
corresponds to sampling a uniformly random branch at each node.
Any capture by the two-point walk gives a capture in $\mathcal{T}(m)$
(\lean{EM/Stochastic/MixedMC.lean}{two\_point\_capture\_implies\_mixed\_capture}
{\color{leanink}$\checkmark$}).

\paragraph{Basic properties.}
Under validity and $m \geq 2$, the formalization proves:
every factor is prime
(\lean{EM/Stochastic/MixedWalk.lean}{mixedWalkFactor\_prime}
{\color{leanink}$\checkmark$});
every factor divides $P{+}1$
(\lean{EM/Stochastic/MixedWalk.lean}{mixedWalkFactor\_dvd}
{\color{leanink}$\checkmark$});
the accumulator stays $\geq 2$
(\lean{EM/Stochastic/MixedWalk.lean}{mixedWalkProd\_ge\_two}
{\color{leanink}$\checkmark$});
the accumulator is strictly increasing
(\lean{EM/Stochastic/MixedWalk.lean}{mixedWalkProd\_strict\_mono}
{\color{leanink}$\checkmark$});
each factor is coprime to the accumulator
(\lean{EM/Stochastic/MixedWalk.lean}{mixedWalkFactor\_not\_dvd\_prod}
{\color{leanink}$\checkmark$});
prefix dependence: $P_\sigma(k)$ depends only on $\sigma(0),\ldots,\sigma(k{-}1)$
(\lean{EM/Stochastic/MixedWalk.lean}{mixedWalkProd\_depends\_on\_prefix}
{\color{leanink}$\checkmark$});
tail restart: the subtree from step~$K$ is a fresh tree from~$P_\sigma(K)$
(\lean{EM/Stochastic/MixedWalk.lean}{mixedWalkProd\_tail\_restart}
{\color{leanink}$\checkmark$}).

\begin{theorem}[Mixed walk structure]
\label{thm:mixed-walk-structure}
For any valid selection~$\sigma$ from a starting point $m \geq 2$:
\begin{enumerate}[nosep]
\item \textbf{Exponential growth}
  (\lean{EM/Stochastic/TreeSieveDecay.lean}{mixedWalkProd\_exp\_growth}
  {\color{leanink}$\checkmark$}):
  $P_\sigma(k) \geq m \cdot 2^k$ for all~$k$.
\item \textbf{Squarefree preservation}
  (\lean{EM/Stochastic/TreeSieveDecay.lean}{mixedWalkProd\_squarefree}
  {\color{leanink}$\checkmark$}):
  if $m$ is squarefree, so is $P_\sigma(k)$ for all~$k$.
\item \textbf{Coprimality from no death}
  (\lean{EM/Stochastic/TreeSieveDecay.lean}{mixedWalkProd\_coprime\_of\_no\_death}
  {\color{leanink}$\checkmark$}):
  if the walk never hits $-1 \bmod q$, then $\gcd(P_\sigma(k), q) = 1$ for all~$k$.
\item \textbf{Strict monotonicity}
  (\lean{EM/Stochastic/MixedWalk.lean}{mixedWalkProd\_strict\_mono}
  {\color{leanink}$\checkmark$}):
  $P_\sigma$ is strictly increasing.
\end{enumerate}
\end{theorem}

\paragraph{Unconditional capture of $q = 3$.}
From $m = 2$, the root has $P{+}1 = 3$.  Since $3$ is prime, every
branch at depth~1 selects factor~$3$---there is no other choice.
Thus $3$ appears in $\mathcal{T}(2)$ unconditionally
(\lean{EM/Stochastic/MixedWalk.lean}{mixed\_capture\_three}
{\color{leanink}$\checkmark$}).

\paragraph{The diversity hypothesis.}
The tree has branching $> 1$ at a node when $P{+}1$ is composite.
\defname{MixedDiversity}
(\lean{EM/Stochastic/MixedWalk.lean}{MixedDiversity}):
along the leftmost path, $P(n){+}1$ is composite cofinally often.
Under MixedDiversity, the tree has genuine branching at infinitely
many levels, and PRE (Theorem~\ref{thm:pre}) guarantees the factors
generate the full unit group $(\ZZ/q\ZZ)^\times$.
Since $P(n)$ grows super-exponentially, $P(n){+}1$ is composite for
all known $n \geq 1$, so MixedDiversity is empirically unchallenged.
The conditional reduction gives:

\begin{theorem}[Diversity implies MixedMC ---
\lean{EM/Stochastic/MixedMC.lean}{mixed\_diversity\_weak\_implies\_mixed\_mc}
{\color{leanink}$\checkmark$}]
\label{thm:diversity-mixed-mc}
$\mathrm{MixedDiversityWeak} + \mathrm{MixedDiversity}
\;\Longrightarrow\; \mathrm{MixedMullinConjecture}$.
The proof is by strong induction on~$q$: MixedDiversityWeak provides
branching at each level, MixedDiversity ensures enough branching
to generate the full unit group, and PRE converts generation to hitting.
\end{theorem}

\paragraph{The perpetual primality obstruction.}
If $P(n){+}1$ were prime for all $n \geq N$, the recurrence
would collapse to $P(n{+}1)+1 = \Phi_3(P(n))$ (the third cyclotomic
polynomial,
\lean{EM/Stochastic/PerpetualPrimality.lean}{perpetual\_prime\_cyclotomic}
{\color{leanink}$\checkmark$}),
and $P(n) \equiv 1 \pmod{3}$ would force $3 \mid \Phi_3(P(n))$,
making $P(n){+}1$ composite---a contradiction.  More precisely,
perpetual primality confines the walk to $\{0,2\} \bmod 3$
(\lean{EM/Stochastic/PerpetualPrimality.lean}{perpetual\_prime\_excludes\_mod3\_one}
{\color{leanink}$\checkmark$}).
More generally, under perpetual primality the walk mod~$q$ follows
the autonomous map $w \mapsto w(w{+}1)$ on $\ZZ/q\ZZ$
(\lean{EM/Stochastic/PerpetualPrimality.lean}{perpetual\_prime\_autonomous\_mod}
{\color{leanink}$\checkmark$}),
and by pigeonhole is eventually periodic
(\lean{EM/Stochastic/PerpetualPrimality.lean}{perpetual\_prime\_eventually\_periodic}
{\color{leanink}$\checkmark$}).
For $q = 5$, if the walk enters~$\{1,2\}$, the autonomous orbit
alternates between $1$ and $2$ forever (since $1 \cdot 2 = 2$ and
$2 \cdot 3 = 1$ in~$\ZZ/5\ZZ$), never reaching $-1 = 4$
(\lean{EM/Stochastic/PerpetualPrimality.lean}{perpetual\_prime\_mod5\_orbit}
{\color{leanink}$\checkmark$}).
These results combine into a clean structural statement:

\begin{theorem}[Perpetual primality periodicity]
\label{thm:perpetual-periodicity}
Suppose $P(n)+1$ is prime for all $n \geq N$.  Then for every prime~$q$:
\begin{enumerate}[nosep]
\item The walk mod~$q$ satisfies the autonomous recurrence
  $w(n{+}1) = w(n) \cdot (w(n){+}1)$
  (\lean{EM/Stochastic/PerpetualPrimality.lean}{perpetual\_prime\_autonomous\_mod}
  {\color{leanink}$\checkmark$}).
\item The walk is eventually periodic mod~$q$ with period $T \leq q$
  (\lean{EM/Stochastic/PerpetualPrimality.lean}{perpetual\_prime\_eventually\_periodic}
  {\color{leanink}$\checkmark$}).
\item For $q = 3$: the walk avoids $1 \bmod 3$, confined to $\{0, 2\}$
  (\lean{EM/Stochastic/PerpetualPrimality.lean}{perpetual\_prime\_excludes\_mod3\_one}
  {\color{leanink}$\checkmark$}).
\item For $q = 5$: the orbit from $\{1, 2\}$ is a $2$-cycle, never reaching
  $-1 = 4$
  (\lean{EM/Stochastic/PerpetualPrimality.lean}{perpetual\_prime\_mod5\_orbit}
  {\color{leanink}$\checkmark$}).
\item For $q = 11$: the orbit from $\{2, 6, 9\}$ is a $3$-cycle, never
  reaching $-1 = 10$
  (\lean{EM/Stochastic/PerpetualPrimality.lean}{perpetual\_prime\_mod11\_orbit}
  {\color{leanink}$\checkmark$}).
\item For $q = 17$: the orbit from $\{2, 3, 4, 6, 8, 12\}$ is a closed
  set (tail~$+$ $2$-cycle $\{3, 12\}$), never reaching $-1 = 16$
  (\lean{EM/Stochastic/PerpetualPrimality.lean}{perpetual\_prime\_mod17\_orbit}
  {\color{leanink}$\checkmark$}).
\item For $q = 23$: the orbit from $\{2, 6, 12, 18, 19, 20\}$ is a closed
  set (tail~$+$ $5$-cycle), never reaching $-1 = 22$
  (\lean{EM/Stochastic/PerpetualPrimality.lean}{perpetual\_prime\_mod23\_orbit}
  {\color{leanink}$\checkmark$}).
\end{enumerate}
The four-prime orbit exclusion
(\lean{EM/Stochastic/PerpetualPrimality.lean}{perpetual\_primality\_multi\_exclusion}
{\color{leanink}$\checkmark$})
shows that under perpetual primality, primes $5, 11, 17, 23$ are
simultaneously blocked from entering the sequence.
\end{theorem}

\noindent
These statements are for the \emph{mixed} walk
$\mathrm{mixedWalkProd}$ of this section.  The corresponding results
for the standard walk---where perpetual primality excludes every prime
$q \equiv 2 \pmod 3$, $q \geq 5$, and refutes MC outright---are in
\S\ref{sec:autonomous-branch}.

%% =========================================================================
\paragraph{Factor Confinement}
\label{sec:factor-confinement}
%% =========================================================================

The reachable set $R_\infty(q, m) = \bigcup_n R_n(q, m)$ cannot be
confined to any proper coset $g \cdot H$ of $(\ZZ/q\ZZ)^\times$:
if $R_\infty \subseteq g \cdot H$, then all multiplier residues lie
in~$H$, contradicting PrimeResidueEscape
(\lean{EM/Equidist/Bootstrap.lean}{prime\_residue\_escape}
{\color{leanink}$\checkmark$}).

\begin{theorem}[\lean{EM/Stochastic/ReachableSets.lean}{reachableEver\_not\_in\_coset}
{\color{leanink}$\checkmark$}]
\label{thm:coset-impossibility}
For $q \geq 3$ prime and any proper subgroup $H < (\ZZ/q\ZZ)^\times$,
$R_\infty(q, 2) \not\subseteq g \cdot H$ for every~$g$.
\end{theorem}

\noindent
Coset impossibility constrains the geometry of $R_\infty$ but does
not force $R_\infty = (\ZZ/q\ZZ)^\times$.  A complementary approach
attacks the \emph{factor} side: if $R_\infty \subsetneq
(\ZZ/q\ZZ)^\times$, then the prime factors of every Euclid number
$P(n){+}1$ along \emph{every} path are confined to a ``forbidden''
set determined by~$R_\infty$
(\lean{EM/Stochastic/ReachableSets.lean}{factor\_confinement}
{\color{leanink}$\checkmark$}).

\begin{property}[FactorEscapeHypothesis ---
\lean{EM/Stochastic/ReachableSets.lean}{FactorEscapeHypothesis}]
\label{prop:factor-escape}
For every prime $q \geq 3$, the Euclid numbers of $\mathcal{T}(2)$
eventually have a prime factor outside the forbidden set induced by any
proper $R \subsetneq (\ZZ/q\ZZ)^\times$.
\end{property}

\begin{theorem}[\lean{EM/Stochastic/ReachableSets.lean}{factor\_escape\_implies\_mixed\_mc\_at}
{\color{leanink}$\checkmark$}]
$\mathrm{FactorEscapeHypothesis}(q) \;\Longrightarrow\;
R_\infty(q, 2) = (\ZZ/q\ZZ)^\times
\;\Longrightarrow\;$ every prime~$q$ appears in $\mathcal{T}(2)$.
\end{theorem}

\noindent
The comparison between the standard walk and the factor tree is stark:
\smallskip
\begin{center}
\begin{tabular}{@{}l@{\quad}c@{\quad}c@{}}
\toprule
 & \textbf{Standard walk} & \textbf{Factor tree} $\mathcal{T}(m)$ \\
\midrule
Factors used per step & $1$ ($\minFac$) & All $\omega(P{+}1)$ \\
Confinement forces & $\minFac \in F$ & Every factor $\in F$ \\
Population analogy & CME (orbit) & PSCD (sieve theory) \\
\bottomrule
\end{tabular}
\end{center}

\smallskip\noindent
For the standard walk, avoiding~$-1$ confines only one factor per
step; for the tree, it confines \emph{all} $\omega(P{+}1)$ factors
simultaneously---a much stronger sieve constraint that can be
attacked with population methods.

%% =========================================================================
\subsection{Almost All GenMixedMC: A Density-1 Result}
\label{sec:almost-all-mixed}
%% =========================================================================

The factor confinement of \S\ref{sec:factor-confinement} connects
the hitting question to sieve theory over the population of starting
points.  For a squarefree $m$ coprime to~$q$, if
$-1 \notin R_\infty(q,m)$, then $R_\infty$ misses a
\emph{nonzero} element of $\ZZ/q\ZZ$.  Moreover, the walk stays in
units (each multiplier is prime and $\neq q$, since $q \mid P{+}1$
would place~$-1$ in~$R_\infty$), so
$0 \notin R_\infty$ as well
(\lean{EM/Ensemble/MixedEnsemble.lean}{zero\_not\_reachable\_of\_coprime\_trapped}
{\color{leanink}$\checkmark$}).
At the root, all prime factors of $m{+}1$ have residues in the
allowed factor set determined by~$R_\infty$
(\lean{EM/Ensemble/MixedEnsemble.lean}{unconditional\_confinement}
{\color{leanink}$\checkmark$}).
By pigeonhole over the finitely many proper finsets~$R$, the
``trapped'' count is bounded by the sum of confined counts
(\lean{EM/Ensemble/MixedEnsemble.lean}{trapped\_le\_sum\_confined}
{\color{leanink}$\checkmark$}).

The pivotal statement \defname{PSCD}
(\emph{Population Sieve Confinement Decay}) --- an open hypothesis
until the chain below closed it --- asserts that among
squarefree integers, the fraction confined to any finset~$R$ missing
a nonzero element of $\ZZ/q\ZZ$ tends to zero.
The restriction to $R$ missing a
\emph{nonzero} element is essential---$R =
(\ZZ/q\ZZ)^\times$ is proper as a subset of $\ZZ/q\ZZ$
but confinement to it merely means $q \nmid m{+}1$,
which has density $(q{-}1)/q > 0$.

\begin{theorem}[\lean{EM/Ensemble/MixedEnsemble.lean}{pscd\_implies\_almost\_all\_mixed\_hitting}
{\color{leanink}$\checkmark$}]
\label{thm:almost-all-mixed}
$\mathrm{PSCD}(q)$ implies: for a density-$1$ set of squarefree~$m$
coprime to~$q$, every prime~$q$ appears in $\mathcal{T}(m)$.
\end{theorem}

\paragraph{Reduction to analytic number theory.}
PSCD reduces to standard analytic number theory via two routes.

\smallskip\noindent
\emph{Route~1 (original, one open bridge).}\enspace
Three intermediate steps:
\begin{enumerate}[nosep]
\item {\color{leanink}$\checkmark$}\enspace
  PrimesEquidistributedInAP $\Rightarrow$ $\sum_{p \equiv a} 1/p$ diverges
  for every coprime class~$a$
  (\lean{EM/Ensemble/MixedEnsemble.lean}{peap\_implies\_fcd\_proved}).
\item \defname{SieveUpperBound}: confined count $\leq$ $C \cdot
  \text{sqfreeCount} \cdot \prod_{p \in F} (1 - 1/p)$
  (fundamental lemma of sieve).
\item {\color{leanink}$\checkmark$}\enspace
  Divergent $\sum 1/p$ implies sieve product $\to 0$
  (\archived{EM/Archive/Ensemble/MixedEnsembleArchive.lean}{sieve\_product\_vanishing\_proved}).
\end{enumerate}
SieveUpperBound---the fundamental lemma of sieve theory---was never
proved here and is now archived: Route~2 made it unnecessary.

\smallskip\noindent
\emph{Route~2 (WeakFMCD, unconditional).}\enspace
The key insight is that FMCD's exact constant is unnecessary:
the sieve product $\prod_{p \in F}(1-1/p) \to 0$ absorbs any
finite multiplicative constant.  Define \defname{WeakFMCD}
(\lean{EM/Ensemble/MixedEnsemble.lean}{weak\_fmcd\_proved}
{\color{leanink}$\checkmark$}):
for any fixed finite set of primes~$S$,
\[
  \operatorname{sqfreeSievedCount}(S,X) \;\leq\;
  \Bigl(4\prod_{p \in S}(1 - 1/p) + \varepsilon\Bigr)
  \cdot \operatorname{sqfreeCount}(X).
\]
The constant~$4$ comes from the unconditional bound
$\operatorname{sqfreeCount}(X) \geq X/4$
(\lean{EM/Ensemble/FirstMoment.lean}{sqfreeCount\_ge\_quarter\_real}
{\color{leanink}$\checkmark$}), proved via $\sum_{d=2}^{D}1/d^2 < 3/4$.
WeakFMCD is then immediate from CRT residue class counting
(each class has $\leq X/N + 2$ elements) and Euler's product
$\varphi(\prod S)/\prod S = \prod(1-1/p)$.

The central theorem
(\lean{EM/Ensemble/MixedEnsemble.lean}{weak\_fmcd\_fcd\_implies\_pscd}
{\color{leanink}$\checkmark$})
proves ForbiddenClassDivergent alone $\Rightarrow$ PSCD\@.  The proof
decomposes the confined count over residue classes mod~$q$ as before,
but with $\varepsilon/(8q)$ for the sieve target and $\varepsilon/(2q)$
for the WeakFMCD tolerance, so $4 \cdot \varepsilon/(8q) + \varepsilon/(2q)
= \varepsilon/q$; summing over $q{-}1$ classes gives $< \varepsilon$.
\paragraph{Closing the chain: Dirichlet density, unconditionally.}
The chain above consumed one analytic input --- divergence of
$\sum_{p \equiv a\,(q)} 1/p$ for every invertible class~$a$
(ForbiddenClassDivergent).  This is now a \emph{theorem} of the
formalization
(\lean{EM/IK/DirichletDensity.lean}{prime\_reciprocal\_class\_divergent}):
the classical Dirichlet-density argument, machine-checked.  The proof
splits the Euler logarithm $\sum_p -\log(1 - \chi(p)p^{-\sigma})$ into
the prime sum plus a tail bounded uniformly by $\sum_p p^{-2}$ for
$\sigma \geq 1$; character orthogonality isolates the class~$a$; the
principal character contributes $\sum_p p^{-\sigma} \to \infty$ as
$\sigma \to 1^+$ (from the divergence of $\sum 1/p$); and for
$\chi \neq \chi_0$ the prime sum stays bounded on $\sigma \in (1,2]$
because $L(\cdot,\chi)$ is continuous and nonvanishing on $[1,2]$ ---
the deep input is $L(1,\chi) \neq 0$ --- so its logarithm along the
segment, tracked through the covering map $\exp$ (path lifting, no
hand-rolled branch bookkeeping), is bounded by compactness.  Mathlib
supplies the $L$-function nonvanishing; the divergence statement
itself fills a gap in Mathlib, which has only the
$\Lambda$-weighted analogue.

\begin{keyresult}
The composed chain
(\lean{EM/Ensemble/UnconditionalPSCD.lean}{almost\_all\_mixed\_hitting\_unconditional}
{\color{leanink}$\checkmark$}):
\[
  \varnothing
  \;\xRightarrow{\text{\color{leanink}$\checkmark$}}\;
  \text{ForbiddenClassDivergent}
  \;\xRightarrow{\text{\color{leanink}$\checkmark$}}\;
  \text{PSCD}(q)
  \;\xRightarrow{\text{\color{leanink}$\checkmark$}}\;
  \text{a.a.\ tree hitting of }q.
\]
Every arrow is machine-checked and the leftmost node is empty:
\textbf{the chain carries no hypothesis at all}.  This is the
formalization's first fully closed nontrivial chain; its axiom
footprint is exactly Lean's three standard axioms
(\code{propext}, \code{Classical.choice}, \code{Quot.sound}).
\end{keyresult}

\paragraph{The diagonal strengthening: all small primes simultaneously.}
The chain above is per-$q$: the density-$1$ set of good starting
points depends on the target prime, and countably many density-$1$
sets need not intersect in a density-$1$ set, so ``almost all~$m$
satisfy GenMixedMC'' does \emph{not} follow.  The strongest
quantifier arrangement derivable from the per-$q$ limits without
uniform-in-$q$ sieve rates is the \emph{diagonal}: there is a bound
$B(X) \to \infty$ such that the density of squarefree $m \leq X$
trapped for \emph{some} prime $q \leq B(X)$ tends to zero
(\lean{EM/Ensemble/UnconditionalPSCD.lean}{almost\_all\_mixed\_hitting\_diagonal}
{\color{leanink}$\checkmark$}) --- almost all starting points reach
every prime up to $B(X)$ \emph{simultaneously}.  The same bound
transfers to the stochastic variant: almost all squarefree starting
points capture, under the $(1{-}\varepsilon)$ process, every prime up
to $B(X)$ simultaneously, for every $\varepsilon > 0$
(\lean{EM/Stochastic/PositiveProbCapture.lean}{almost\_all\_positive\_prob\_capture\_diagonal}
{\color{leanink}$\checkmark$}).  Density-one GenMixedMC itself ---
all primes at once from a single density-$1$ set --- would require
uniform-in-$q$ rates in the sieve and remains open.

\paragraph{Missed primes: forward invariance and the $q=3$ model case.}
For a starting point~$m$ write $\mathrm{Miss}(m)$ for its
\emph{missed-prime set}: the primes~$q$ coprime to~$m$ with
$-1 \notin R_\infty(q,m)$.  Missing is \emph{forward-invariant} along
the factor tree: $\mathrm{Miss}(m) \subseteq \mathrm{Miss}(mf)$ for
every child~$mf$ with $f$ prime, $f \mid m{+}1$
(\lean{EM/Stochastic/MissedPrimes.lean}{missed\_primes\_forward\_invariant}
{\color{leanink}$\checkmark$}) --- equivalently, if any tree node
reaches~$q$ then so does the root.  Conversely, when $-1$ is reachable
but not yet attained, some child still reaches it
(\lean{EM/Stochastic/MissedPrimes.lean}{good\_child\_exists}
{\color{leanink}$\checkmark$}) --- the stepwise invariant any
almost-sure argument must maintain.  For $q = 3$ the reachability side
closes \emph{pointwise}: $3 \notin \mathrm{Miss}(m)$ for every
$m \geq 2$
(\lean{EM/Stochastic/MissedPrimes.lean}{three\_notMem\_missedPrimes}
{\color{leanink}$\checkmark$}), and capture opportunities for~$3$ are
\emph{cofinal}: from every squarefree $m \geq 2$ coprime to~$3$ and
beyond every horizon~$K$, some valid path has $3 \mid P_n{+}1$ with
$n \geq K$ --- unless capture is forced earlier because $P_n{+}1$ is a
power of~$3$
(\lean{EM/Stochastic/MissedPrimes.lean}{three\_cofinal\_capture\_opportunities}
{\color{leanink}$\checkmark$}); in particular positive-probability
capture of~$3$ holds pointwise, from \emph{every} such start
(\lean{EM/Stochastic/MissedPrimes.lean}{three\_positive\_prob\_capture\_pointwise}
{\color{leanink}$\checkmark$}).  On the ensemble side, the expected
number of missed primes $q \leq B(X)$ of a random squarefree start,
with $B(X) \to \infty$, tends to~$0$
(\lean{EM/Ensemble/UnconditionalPSCD.lean}{expected\_missed\_smallprimes\_diagonal}
{\color{leanink}$\checkmark$}).  What separates the above from
\emph{almost-sure} capture of~$3$ under the $(1{-}\varepsilon)$
process is pure anatomy, not reachability: along a path with cofinal
opportunities the process fails them all with probability at most
$\prod_k (1 - \varepsilon/\omega(P_k{+}1))$, which vanishes iff
$\sum_k 1/\omega(P_k{+}1) = \infty$ --- so almost-sure capture of~$3$
needs only an $\omega$-bound on Euclid numbers along walks (true for
the normal order $\omega \sim \log\log P_k \sim k$, unprovable from
the trivial $\omega \leq \log_2 P_k$), plus a path-measure layer.
Regeneration, the open gap for general~$q$, is thus \emph{free} for
$q = 3$.

\begin{deadend}[title={Dead End \#117: MultCancelToWalkCancel}]
The standard ensemble approach
(\S\ref{sec:ensemble-variant}) proves almost-all
\emph{multiplier cancellation} for the deterministic walk, but
transferring this to walk \emph{hitting} is blocked by
Dead End~\#117 (MultCancelToWalkCancel).  The factor tree approach
bypasses~\#117 entirely:
\smallskip
\begin{center}
\begin{tabular}{@{}l@{\quad}l@{\quad}l@{}}
\toprule
 & \textbf{Standard ensemble} & \textbf{Factor tree ensemble} \\
\midrule
Population result & a.a.\ char cancel & a.a.\ GenMixedMC \\
Transfer gap & \#117 (cancel $\not\Rightarrow$ hitting) & None \\
Selection used & $\minFac$ only & All factors of $P{+}1$ \\
Open hypothesis & centered StepDecorrelation (uncentered form false, \#156) & None (chain closed) \\
\bottomrule
\end{tabular}
\end{center}
\end{deadend}

\smallskip\noindent
This is the strongest density result in the formalization, and it is
\emph{unconditional}: for each prime~$q$, a density-$1$ set of
squarefree starting points reaches~$q$ in the factor tree (the good
set depends on~$q$; see the diagonal strengthening below).

\paragraph{Interpolation: tree reachability implies positive probability capture.}
The density-$1$ tree result bridges to the $(1{-}\varepsilon)\,\minFac
+ \varepsilon\,$random process.  Assign to each factor choice at
accumulator~$P$ the weight $\varepsilon / \omega(P{+}1)$---a lower
bound on the probability of that choice under the interpolated process
(where $\omega$ counts distinct prime factors).  A finite path
$\sigma$ of length~$n$ has weight
$\prod_{k<n} \varepsilon / \omega(P_\sigma(k){+}1) > 0$
(\lean{EM/Stochastic/PositiveProbCapture.lean}{pathWeightLB\_pos}
{\color{leanink}$\checkmark$}).

If $-1 \in R_\infty(q,m)$ (tree reachability), a capturing path
exists, and that path has positive weight under the
$(1{-}\varepsilon)$ process for \emph{any} $\varepsilon > 0$
(\lean{EM/Stochastic/PositiveProbCapture.lean}{reachable\_implies\_positive\_prob\_capture}
{\color{leanink}$\checkmark$}).
Almost all squarefree~$m$ have $-1$ reachable \emph{unconditionally}
(the Dirichlet-density chain of \S\ref{sec:almost-all-mixed}), so
almost all have positive-probability capture with \textbf{no open
hypothesis}
(\lean{EM/Stochastic/PositiveProbCapture.lean}{almost\_all\_positive\_prob\_capture\_unconditional}
{\color{leanink}$\checkmark$}): from almost every squarefree starting
point, the $(1{-}\varepsilon)$ process can reach any given prime, for
every $\varepsilon > 0$.  As with every density statement here, the
good set depends on the target prime~$q$; any \emph{finite} set of
primes shares a common density-$1$ set of starting points.  (The
original PEAP-conditional route is retained as
\lean{EM/Stochastic/PositiveProbCapture.lean}{almost\_all\_positive\_prob\_capture}.)

\paragraph{Layer 2: cofinal and iterated hitting via regeneration.}
Layer~1 proves a \emph{single} hit with positive weight.  Stronger
conclusions require \defname{Regeneration}
(\lean{EM/Stochastic/TreeSieveDecay.lean}{Regeneration}): every
accumulator encountered along any valid walk from~$m$ is \emph{good},
i.e., every unit of~$(\ZZ/q\ZZ)^\times$ remains reachable.  The
evidence is strong: accumulators grow super-exponentially and are
squarefree, so by CRT blindness they should behave generically mod~$q$.

Under Regeneration, the formalization proves:
\begin{itemize}
\item \textbf{Block coverage}
(\lean{EM/Stochastic/TreeSieveDecay.lean}{block\_coverage}
{\color{leanink}$\checkmark$}):
from a good accumulator, there is a uniform depth bound~$N_0$ such
that every unit can be reached within~$N_0$ steps.
\item \textbf{Cofinal hitting}
(\lean{EM/Stochastic/TreeSieveDecay.lean}{regeneration\_implies\_cofinal\_hitting}
{\color{leanink}$\checkmark$}):
for every threshold~$K$, there exists a valid path that hits
$-1\bmod q$ at some step~$n \geq K$, with positive weight.
\item \textbf{Iterated hitting}
(\lean{EM/Stochastic/TreeSieveDecay.lean}{regeneration\_implies\_iterated\_hitting}
{\color{leanink}$\checkmark$}):
for every target number of $-1$ hits, there exists a valid path
achieving at least that many hits, each with positive weight.
The proof proceeds by induction on the target, concatenating walks
via \code{concatMixedSelection} and using Regeneration to obtain
fresh good accumulators past each existing hit.
\end{itemize}
Regeneration is the sole open hypothesis in this layer; it is
strictly weaker than CME (which implies MC outright) and concerns
the persistence of a qualitative property (full reachability) along
the walk, not a quantitative equidistribution statement.

\paragraph{TreeSieveDecay: the sieve-theoretic bridge.}
Regeneration asks that \emph{every} intermediate accumulator is good.
A weaker, sieve-flavored variant is \defname{TreeSieveDecay}
(\lean{EM/Stochastic/TreeSieveDecay.lean}{TreeSieveDecay}):
there exists a threshold~$P_0$ such that every squarefree~$P \geq P_0$
coprime to~$q$ is a good accumulator.
TSD is a ``large-$P$'' hypothesis: all sufficiently generic starting
points support full reachability.
The coprimality condition is necessary: when $q \mid P$, absorption
traps the walk at~$0$ permanently, so no unit is reachable.
TSD bypasses Regeneration entirely---for a single hit (which suffices
for MC), a cleaner dichotomy argument works.
The formalization proves two bridges:
\begin{itemize}
\item \textbf{Squarefree propagation}
(\lean{EM/Stochastic/TreeSieveDecay.lean}{mixedWalkProd\_squarefree}
{\color{leanink}$\checkmark$}):
if the starting point is squarefree, every intermediate accumulator
is squarefree (induction, using coprimality of factors).
\item \textbf{Coprimality from no death}
(\lean{EM/Stochastic/TreeSieveDecay.lean}{mixedWalkProd\_coprime\_of\_no\_death}
{\color{leanink}$\checkmark$}):
if the walk never hits $-1 \bmod q$, all intermediate accumulators
stay coprime to~$q$ (if the factor equaled~$q$, then
$q \mid P{+}1$, so the walk position would be~$-1$---contradiction).
\end{itemize}
Together: under TSD, from any sufficiently large squarefree~$m$,
the walk either hits~$-1$ (done) or stays coprime to~$q$, in which
case TSD applies to each intermediate accumulator directly
(\lean{EM/Stochastic/TreeSieveDecay.lean}{tsd\_implies\_neg\_one\_reachable}
{\color{leanink}$\checkmark$}).

A weaker variant \defname{TreeSieveDecayHitting}
(\lean{EM/Stochastic/TreeSieveDecay.lean}{TreeSieveDecayHitting})
requires only that $-1$ is reachable (not all units).
TSD implies TSDH
(\lean{EM/Stochastic/TreeSieveDecay.lean}{tsd\_implies\_tsdh}
{\color{leanink}$\checkmark$}).

\paragraph{Orbit melting: backward propagation.}
Two squarefree positive integers with the same prime factor set are
equal
(\lean{EM/Stochastic/TreeSieveDecay.lean}{eq\_of\_same\_primeFactors\_squarefree}
{\color{leanink}$\checkmark$}).
This ``orbit melting'' identity means the accumulator carries no
hidden state beyond its prime factorization:
same accumulator $\Rightarrow$ same reachable set
(\lean{EM/Stochastic/TreeSieveDecay.lean}{same\_accumulator\_same\_future}
{\color{leanink}$\checkmark$}).
More importantly, a good accumulator at \emph{any} intermediate step
propagates backward: if $P_\sigma(K)$ supports full reachability,
then $-1$ is reachable from the starting point~$m$ by concatenating
the prefix path with the witnessing path
(\lean{EM/Stochastic/TreeSieveDecay.lean}{good\_accumulator\_propagates\_to\_start}
{\color{leanink}$\checkmark$}).

\paragraph{TSD-Hitting(3): unconditional.}
For $q = 3$, TreeSieveDecayHitting holds \textbf{unconditionally}---no
open hypothesis at all.  The argument is a mod-$3$ dichotomy on any
squarefree~$P$ coprime to~$3$:
\begin{itemize}
\item If $P \equiv 2 \pmod{3}$: then $(P : \ZZ/3\ZZ) = -1$, so the
  walk is already at~$-1$ at step~$0$.
\item If $P \equiv 1 \pmod{3}$: then $P{+}1 \equiv 2 \pmod{3}$.
  A number $\geq 2$ that is $\equiv 2 \bmod 3$ and not divisible by~$3$
  must have a prime factor $p \equiv 2 \pmod{3}$
  (\lean{EM/Stochastic/TreeSieveDecay.lean}{exists\_factor\_of\_succ\_mod3}
  {\color{leanink}$\checkmark$},
  by strong induction: if $\minFac \equiv 1$, divide out and recurse).
  Selecting that factor gives $P \cdot p \equiv 1 \cdot 2 = -1 \pmod{3}$
  at step~$1$.
\end{itemize}
\begin{keyresult}
\begin{theorem}[\lean{EM/Stochastic/TreeSieveDecay.lean}{tsd\_hitting\_three\_unconditional}]
$\mathrm{TreeSieveDecayHitting}(3)$ holds unconditionally: for every
squarefree $P \geq 2$ coprime to~$3$, there exists a valid path in
the factor tree that reaches $-1 \bmod 3$ within one step.
\end{theorem}
\end{keyresult}

\begin{leancode}
-- The hypothesis: -1 reachable from all large squarefree coprime P
def TreeSieveDecayHitting (q : Nat) : Prop :=
  exists P_0, forall P, P_0 <= P -> Squarefree P -> Nat.Coprime P q ->
    (-1 : ZMod q) in reachableEver q P

-- UNCONDITIONAL for q = 3 (zero open hypotheses)
theorem tsd_hitting_three_unconditional : TreeSieveDecayHitting 3 := by
  -- mod-3 dichotomy: P == 2 -> already at -1;
  -- P == 1 -> P+1 == 2, so some prime factor p == 2 mod 3 exists
  ...
\end{leancode}

\noindent
This is the first unconditional \emph{tree-level} hitting result in this
formalization beyond the trivial $q = 3$ capture of
\S\ref{sec:factor-confinement}.
The previous $q{=}3$ result
(\lean{EM/Stochastic/MixedWalk.lean}{mixed\_capture\_three}
{\color{leanink}$\checkmark$})
showed that $3$ appears in $\mathcal{T}(2)$ from the fixed root~$m{=}2$;
TSD-Hitting(3) proves that $-1$ is reachable from \emph{every}
coprime squarefree accumulator---a uniform statement that applies at
any intermediate point of the walk.

\paragraph{TSD $\Rightarrow$ MixedMC: composition.}
The threshold condition in TSD (accumulator $P$ must exceed
a bound~$P_0$) is eliminated by the exponential growth of the mixed
walk: the accumulator $P_\sigma(k)$ satisfies
$P_\sigma(k) \geq m \cdot 2^k$
(\lean{EM/Stochastic/TreeSieveDecay.lean}{mixedWalkProd\_exp\_growth}
{\color{leanink}$\checkmark$}),
since each multiplier is a prime $\geq 2$.  For any starting
accumulator $m \geq 2$, the walk eventually exceeds~$P_0$; a
death/no-death dichotomy then applies: either the walk already hit
$-1$ (done), or all accumulators remain coprime to~$q$ and
squarefree, so TSD applies at the first step past the threshold, and
tail reachability (Part~9) propagates $-1$ back to the root.  This
gives the full chain:

\begin{theorem}[TSD reduction chain]
\label{thm:tsd-chain}
For each prime $q \geq 5$:
\[
  \mathrm{TSD}(q)
  \;\Longrightarrow\;
  \forall\, m \geq 2 \text{ sqfree, } \gcd(m,q)=1 \colon
  {-}1 \in R_\infty(q, m)
  \;\Longrightarrow\;
  \mathrm{MixedMC}(q).
\]
The first arrow uses exponential growth
(Theorem~\textup{\ref{thm:mixed-walk-structure}})
to exceed the TSD threshold
(\lean{EM/Stochastic/TreeSieveDecay.lean}{tsd\_implies\_reachable\_from\_any}
{\color{leanink}$\checkmark$});
the second uses hit-implies-capture
(\lean{EM/Stochastic/MixedMC.lean}{hit\_implies\_capture'}
{\color{leanink}$\checkmark$}).
Consequently
(\lean{EM/Stochastic/TreeSieveDecay.lean}{tsd\_all\_implies\_mixed\_mc}
{\color{leanink}$\checkmark$}):
\[
  (\forall q \geq 5,\; \mathrm{TSD}(q))
  \;\Longrightarrow\; \mathrm{MixedMullinConjecture}.
\]
\end{theorem}

\noindent Small primes $q \in \{2, 3\}$ are handled by their existing
unconditional results (\code{mixed\_mc\_two}, \code{mixed\_mc\_three}).

\paragraph{Coset ambiguity: why TSD-Hitting(3) is free but $q \geq 5$ is not.}
The structural reason that TSD-Hitting($q$) closes unconditionally
for $q = 3$ but requires a hypothesis for $q \geq 5$ is \emph{coset
ambiguity}.  For $q = 3$, the group $(\ZZ/3\ZZ)^\times = \{1, 2\}$
has a single non-identity element, which \emph{is} $-1 = 2$: any
prime factor escaping the trivial coset $\{1\}$ automatically
lands on~$-1$
(\lean{EM/Stochastic/TreeSieveDecay.lean}{single\_coset\_implies\_immediate\_hit}
{\color{leanink}$\checkmark$}).
For $q = 5$, the group $(\ZZ/5\ZZ)^\times = \{1,2,3,4\}$ has the
subgroup $\{1, 4\} = \{1, -1\}$ and two other cosets $\{2\}$, $\{3\}$,
neither equal to~$-1$.  A concrete counterexample:
$P = 2$, $P{+}1 = 3$, $(3 \bmod 5) = 3 \neq 4 = -1$---the only
available factor misses~$-1$
(\lean{EM/Stochastic/TreeSieveDecay.lean}{coset\_ambiguity\_counterexample}
{\color{leanink}$\checkmark$}).

\paragraph{Ensemble factor diversity and the stochastic EM process.}
The ensemble-level structure of the prime factor sets
$\text{primeFactors}(\Prod{n}(k)+1)$ is formalized in
\code{EM/Stochastic/FactorDiversity.lean} (346~lines, zero \code{sorry}).

\begin{theorem}[Spectral contraction from factor diversity ---
\lean{EM/Stochastic/FactorDiversity.lean}{factor\_diversity\_spectral\_contraction}
{\color{leanink}$\checkmark$}]
\label{thm:factor-diversity-contraction}
If the prime factor set of $\Prod{n}(k)+1$ contains $\geq 2$ distinct
residue classes mod~$q$, then the mean character value over the factor
set satisfies $\|\bar\chi(F)\| < 1$ for every nontrivial~$\chi$.
If such diverse steps occur infinitely often with a uniform gap
$\delta > 0$, the averaged character product tends to zero
(\lean{EM/Stochastic/FactorDiversity.lean}{diverse\_steps\_imply\_vanishing}
{\color{leanink}$\checkmark$}).
\end{theorem}

\noindent
The bridge from diversity to stochastic capture
(\code{EM/Stochastic/DiverseStepsToCapture.lean}, 281~lines) shows each
prime factor of $\Prod{2}(k)+1$ produces a distinct element of the
reachable set $R_\infty(q, 2)$
(\lean{EM/Stochastic/DiverseStepsToCapture.lean}{genFactor\_in\_reachableEver}
{\color{leanink}$\checkmark$}).

\begin{theorem}[Integer stochastic MC ---
\lean{EM/Stochastic/StochasticEM.lean}{stochastic\_mc\_of\_tsd}
{\color{leanink}$\checkmark$}]
\label{thm:integer-stochastic-mc}
$\mathrm{TSD}(q) \;\Longrightarrow\; \mathrm{StochasticMC}(q)$ for all
$\varepsilon > 0$, where $\mathrm{StochasticMC}(q)$ is the existence of a
finite capturing path of positive weight (tree reachability), not an
almost-sure statement.  The companion
\lean{EM/Stochastic/StochasticEM.lean}{phase\_transition\_summary}
{\color{leanink}$\checkmark$} records only that character values have
norm~$1$ while the mean of a character over a factor set on which it
takes two values has norm~$<1$; we no longer call this a ``sharp phase
transition''.
\end{theorem}

\paragraph{Probabilistic infrastructure.}
The TSD$\to$MixedMC chain is supported by three infrastructure
modules.  \code{EM/Stochastic/TransitionKernel.lean} defines the
epsilon-mixed walk's step PMF and proves the weight lower bound
$\varepsilon/\omega(P{+}1)$ bridges to InterpolationMC's
\code{stepWeightLB}.  \code{EM/Stochastic/PathMeasure.lean} converts
set-theoretic reachable sets to computable Finsets and formalizes
branching distinctness: when $P$ is a unit mod~$q$ and $P{+}1$ has
distinct factor residues, each gives a distinct next walk position.
\code{EM/Stochastic/GeometricCapture.lean} provides the block-geometric
engine: per-block positive capture weight decays non-capture
probability as $(1{-}\delta)^K \to 0$, turning TSD into asymptotic
capture certainty.

